% Shared style for all Git & GitHub lectures (same look as the N-body course)
\documentclass[11pt,a4paper]{article}
\usepackage[margin=2.3cm]{geometry}
\usepackage[T1]{fontenc}
\usepackage{lmodern}
\usepackage{upquote}   % straight quotes in monospaced text, so commands can be copied
\usepackage{microtype}
% keep '--' as two real hyphens in monospace text (commands like --version)
\DisableLigatures[-,<,>]{family=tt*}
\usepackage{booktabs,enumitem,graphicx,xcolor}
\usepackage{array}
\usepackage{needspace}
\usepackage{longtable}
\usepackage{float}
\usepackage{tikz}
\usetikzlibrary{arrows.meta,calc,positioning,decorations.pathreplacing,shapes.geometric,fit,backgrounds}
\usepackage[most]{tcolorbox}
\usepackage{listings}
\usepackage{fancyhdr}
\usepackage[colorlinks=true,linkcolor=blue!50!black,urlcolor=blue!50!black,citecolor=blue!50!black]{hyperref}

\setlength{\parindent}{0pt}
\setlength{\parskip}{0.6em}
\setlength{\abovecaptionskip}{4pt}
\usepackage[font=small,labelfont=bf]{caption}
\setlist{itemsep=0.2em,topsep=0.3em}

% ---------- colours ----------
\definecolor{cbg}{HTML}{F6F8FA}
\definecolor{ckw}{HTML}{0B5CAD}
\definecolor{ccm}{HTML}{6A737D}
\definecolor{cst}{HTML}{B5451B}
\definecolor{cpart}{HTML}{1F6FB2}
\definecolor{cheavy}{HTML}{C0392B}
\definecolor{cgreen}{HTML}{2E8B57}

% ---------- terminal listings ----------
% In a terminal session, lines starting with "$" are what YOU type.
% Everything else is what Git prints back.
\lstdefinestyle{term}{
  basicstyle=\ttfamily\small, backgroundcolor=\color{cbg},
  frame=single, rulecolor=\color{black!15}, showstringspaces=false,
  breaklines=true, columns=fullflexible, keepspaces=true,
  aboveskip=0.6em, belowskip=0.6em, xleftmargin=4pt, xrightmargin=4pt,
  literate={\$\ }{{\textcolor{cgreen}{\$\ }}}2,
  upquote=true
}
\lstset{style=term}
\newcommand{\code}[1]{\texttt{\small #1}}

% ---------- boxes ----------
\newtcolorbox{history}{colback=yellow!8,colframe=orange!60!black,
  title={History},fonttitle=\bfseries\small,left=4pt,right=4pt,top=2pt,bottom=2pt}
\newtcolorbox{quickfact}{colback=purple!4,colframe=purple!55!black,
  title={Quick fact},fonttitle=\bfseries\small,left=4pt,right=4pt,top=2pt,bottom=2pt}
\newtcolorbox{keyidea}{colback=blue!4,colframe=blue!50!black,
  title={Key idea},fonttitle=\bfseries\small,left=4pt,right=4pt,top=2pt,bottom=2pt}
\newtcolorbox{pitfall}{colback=red!4,colframe=red!55!black,
  title={Pitfall},fonttitle=\bfseries\small,left=4pt,right=4pt,top=2pt,bottom=2pt}
\newtcolorbox{tryit}{colback=teal!5,colframe=teal!55!black,
  title={Try it},fonttitle=\bfseries\small,left=4pt,right=4pt,top=2pt,bottom=2pt,breakable}
\newtcolorbox{homework}[1]{colback=green!4,colframe=green!40!black,
  title={Homework #1},fonttitle=\bfseries\small,left=4pt,right=4pt,top=2pt,bottom=2pt,
  breakable}
\newtcolorbox{goals}{colback=gray!6,colframe=gray!60!black,
  title={Learning goals},fonttitle=\bfseries\small,left=4pt,right=4pt,top=2pt,bottom=2pt}

% Solutions always start on a fresh page, behind a clear warning.
\newcommand{\solutionspage}{%
  \clearpage
  \begin{tcolorbox}[colback=red!6,colframe=red!60!black,halign=center]
    {\Large\bfseries Solutions}\\[0.3em]
    {\small Stop here if you have not tried the homework yet.\\
    Attempt every task first, run the checker, and only then read on.}
  \end{tcolorbox}
}

% ---------- Git diagram styles (TikZ) ----------
\tikzset{
  commit/.style={circle,draw=black!70,fill=blue!10,minimum size=7mm,inner sep=0pt,font=\scriptsize\ttfamily},
  newcommit/.style={commit,fill=orange!25,dashed},
  tag/.style={draw,rounded corners=2pt,font=\scriptsize\ttfamily,inner sep=2pt},
  branchtag/.style={tag,fill=green!15,draw=green!40!black},
  headtag/.style={tag,fill=yellow!30,draw=orange!70!black},
  remotetag/.style={tag,fill=purple!10,draw=purple!60!black},
  area/.style={draw,rounded corners,minimum height=11mm,minimum width=33mm,align=center,font=\small},
  cmdarrow/.style={->,thick,>=Stealth},
  link/.style={->,thick,>=Stealth,black!70},
}

% ---------- header ----------
\pagestyle{fancy}
\fancyhf{}
\renewcommand{\headrulewidth}{0.3pt}
\fancyfoot[C]{\thepage}
\newcommand{\lecturetitle}[3]{%
  \fancyhead[L]{\small\textit{Git \& GitHub from Zero}}
  \fancyhead[R]{\small Lecture #1}
  \begin{center}
    {\Large\bfseries Lecture #1}\\[0.4em]
    {\LARGE\bfseries #2}\\[0.4em]
    {\small #3}
  \end{center}
  \vspace{0.5em}
}
